\section{Man in the High Castle}

\begin{quotex}
Ostensibly a review of some artifacts of popular culture including the Man in the High Castle, the Vikings and Valkyrie, but a few may notice something else ...
\end{quotex}

Amazon recently released a dramatic television series, \emph{The Man in the High Castle}, loosely based on \textbf{Philip K Dick}'s novel of the same name first published in 1962. The premise is based on the counter-factual history that the Axis powers won World War II. The Pacific states are occupied by Japan and the Eastern states by Germany, with a buffer zone around the Rockies.

It's difficult to understand Amazon's purpose, since the occupied USA seems more attractive. Was that by accident or design? Or most likely, the writers saw it as less attractive. The Nippon states were less attractive, since the white citizens were pretty much in dhimmitude to the Japanese. Nevertheless, one is struck by the Japanese code of honor and politeness, at least among each other. Moreover, a government official who failed in his duties felt shame to the extent that suicide was preferable. Now in the USA, many politicians are willing to ``accept responsibility'' for their failures, but that is something quite different from accepting the consequences.

The Eastern sector was populated by normal families, who were well groomed, polite, dined together, interacted with each other rather than electronic devices, valued patriotism and education. Apparently the intact family, in the mind of the Amazon Studio writers, sounds fascist to the modern mind. On the negative side, the right to bear arms was rejected and Bibles forbidden, available only on the black market. In 1962, a well-bred man in the USA would have been shocked that a government would have that power, since he still ``clung to his gun and his Bible''. Nowadays the possession of guns and Bibles is considered fascist and the utopian left is working to ban both of them; what was fascist in 1962 is progressive today.

If the elimination of guns and bibles sounds normal to you, perhaps assisted suicide still makes you queasy, at least it was unthinkable in 1962. When the son of the American Nazi officer John Smith starts exhibiting symptoms of a genetic degenerative disease, the father is given a suicide kit to use on the son. Nowadays, it is called ``end of life'' counseling, and the ``right to die'' is increasingly accepted. Moreover, the idea of euphemistic post-natal abortions\footnote{\url{http://aleteia.org/2014/11/05/post-birth-abortions-an-idea-thats-gaining-ground-on-campus}} is gaining acceptance. So the idea of a hospital cremating the handicapped and terminally ill, as in the series, is really a foretelling of our own future rather than an imaginary fascist past.

There is even an annoying SJW, Juliana, who is headstrong in taking on a mission she does not understand, oblivious to the collateral damage she is causing. She illustrates the danger of the sex instinct, since her boyfriend, a young nazi secret agent, and even her I Ching casting Japanese boss, all had major crushes on her. Let's just say some people get hurt, especially the boyfriend's sister, niece and nephew. The men all compromise their own life missions for her.

So back to the important part of the story, which we will have to fill in, since it was unclear in the series. The man in the high castle is collecting films that depict alternate versions of reality. For example, in one of them, the Allies won the war, in another the secret agent assassinated the boyfriend.

Now we are getting into Philip K Dick territory, even if he himself did not actually use that motif. In \emph{Gnosis}, \textbf{Boris Mouravieff} compares each person's life to a film in which he has a leading part. Such as we are, we are passive participants in the film, interacting with persons for unknown reasons. However, the goal is to become the conscious agent --- writer and director --- of the script of one's own life.

Perhaps that is the real point of the series, rather than as an homage to our ``superior'' way of life. As Guenon pointed out, ``The end of one world is the beginning of another.'' Awakening from the common dream leads to the real world, but it must start with a few.

\paragraph{Valkyrie}


The series includes a failed assassination attempt on Hitler, perhaps mirroring a real attempt in 1944, that \textbf{Julius Evola}, writing in 1952\footnote{\url{https://gornahoor.wordpress.com/2015/12/06/a-failed-revolution/}}, claimed was little known. \emph{Valkyrie} is the title of a 2008 film, about the 1944 assassination attempt on Hitler, starring Tom Cruise, so it is perhaps more well-known now. One can easily get the impression from these shows that liberals and progressives were interested in eliminating Hitler. However, in fact, the attempt was engineered by the conservative elements in Germany who did not regard Hitler as conservative or right-wing.

\paragraph{The Vikings}


The History Channel has been airing a series about the Vikings at the time they were becoming interested in England and France. The broad outline is historically correct, although certainly not in every detail. There is a lesson in it about a real ``clash of civilizations''. A couple of generations after Charlemagne, Europe had become Christianized while the Vikings remained pagans. Although both sides practiced Realpolitik, the Vikings had no regard for just way theory. Whenever they wanted something, they went on a ``raid''. Saxon England and Paris, as more advanced, and possibly more effete, civilizations, provided attractive targets for raids.

There may be lessons here for contemporary attempts at multi-culturalism: it does not happen by wishful thinking. Even when a Saxon king allowed a Viking settlement on his land, the problems did not go away. Differences in race, religion, language, and culture are not easily bridged.

The conflict between pagans and Christians is one of the more interesting aspects. Despite their general amorality, there is still something noble about the Vikings. They regard patriarchy as normative, support their community, love and value their children, are indifferent to pain, accept the tragic side of life, and look forward to an afterlife. Ragnar, the Viking leader, tells his son that ``happiness does not matter'', i.e., it is neither the purpose of life nor a goal to pursue. So the ``pursuit of happiness'' in the \emph{Declaration of Independence} must be an illusory right.

The Nordic gods and stories are a real presence and even temptation. Nevertheless, from an esoteric perspective all that must be left behind. One Viking explains to a Christian monk that his gods are real, they eat, fight, and so on. To him such humanized gods are regarded as a strength. Yet, \textbf{Boris Mouravieff} tells us something different:

\begin{quotex}
Reason attributes to the divine an attitude, a weakness, and even more often, purely human motives... which tended to humanize the divinities. The Good News announced by Jesus reversed this ancient conception, calling for the divinization of the human in man by a second Birth; the gateway to the Kingdom of God.
\end{quotex}


Some try to idealize the pagan life in its immediacy and closeness to nature. They can't see past the first birth into the state of nature. The task of the esoterist is quite different: it is not to humanize the gods, but rather to divinize the human.

A final thought: in our time, it is rare to see a movie or TV show with an all white cast and no gay characters. If you are flirting with white nationalism, you can look and then decide if you like that. I suspect not, since there is no gay character to explain to the Vikings how to act more manly.

\paragraph{When We Were Young}


Avoid this if you can; I couldn't. It concerns a culture clash between a boomer and an ambitious millennial. There was a minor point that was of interest, viz., acting either from innocence or experience. The boomer was all high tech, but the millennial went retro. For example, he played board games rather than electronic games, used an old mechanical typewriter for his scripts, sported a fedora, and owned an extensive music collection on vinyl LPs.

I don't see that millennials are going retro in droves, although Barnes \& Noble is phasing out CDs in favor of LPs. Now in fact I play board and card games rather than electronic games, often wear a fedora, and still have a nice collection of vinyl. But I do it naively, since that was the world in which I was raised. It is not necessarily a conscious choice, but more a matter of habit.

However, the millennial is acting ironically. He makes a conscious decision to reject the contemporary option; it is not a habit, but rather an acquired taste.

We are in an analogous position. The world of our fathers no longer exists, so we can't absorb that worldview automatically without thinking about it much. Quite the contrary, we are faced with a choice since the worldview of the modern world totally surrounds us. Hence, we can no longer be naïve, but rather we need to know exactly why we adopt one worldview while rejecting its alternatives.

For the innocent man, the will follows the intellect which was formed by family, society, and church. The man of experience needs to consciously create. He is in the situation described by Mouravieff:

\begin{quotex}
By progressively taking his fate into his own hands, man at the same time takes responsibility for all the partners in his film. It has already been said that he must restore the original meaning of his film, then push the development of the latter in such a way that the `play' be properly played out to its intended denouement. The hero, while working on himself, must apply himself to create new circumstances around him, which will enhance the unfolding of the action towards its originally intended conclusion.
\end{quotex}



\flrightit{Posted on 2015-12-17 by Cologero}

\begin{center}* * *\end{center}

\begin{footnotesize}
\begin{sffamily}
\texttt{AghorNath on 2015-12-18 at 18:17 said:}

Thank you for the post. Brilliant and thought provoking as always. I was wondering if ``Man in the High Castle'' was actually worth watching as far as production value, character structure, presentation etc? I remember trying to watch it sometime ago, but didn't give it much of a chance. Once they started showing scenes of NYC decked out in regalia I turned it off. It all looked so plastic and silly, so not what one with any sense of the aesthetics that these powers conveyed would expect. Too overtly `Hollywood', and sensationalized. Any thoughts? Should I give it a go, or have your sentiments surmised and articulated it enough beyond it's scope? Cheers

\hfill

\texttt{Avery on 2015-12-18 at 22:05 said:}

I met a Chinese guy who has adopted an ironic faith in Buddhism. When he sees a Buddha image, he puts his hands together, bows, gets down on both knees and prostrates himself in prayer. Except that he obviously didn't learn this behavior from his parents; he's urban elite and extremely overeducated. The piety comes off as false and pretentious. It makes one want to doubt Guénon's principle that there is no such thing as empty ritual.

\hfill

\texttt{Cologero on 2015-12-18 at 23:32 said:}

@AghorNath, the reviews are not usually intended to promote such popular works. Rather, we are trying to illustrate the third dimension of history by rooting out and exposing the unexamined prejudices of such works. In an intelligent work of art, on the other hand, the artist works consciously, understanding full well what he is doing. He will help the reader (or observer) bring to light his own unexamined prejudices. That is the difference between, for example, the recently reviewed \emph{Submission} and the Amazon production.

\hfill

\texttt{Cologero on 2015-12-18 at 23:36 said:}

@Avery, perhaps he is just doing it badly. Has he performed the tea ceremony for you?

How do you explain \textit{Guenon's Conversion}\footnote{\url{http://www.gornahoor.net/?p=262}}? When he answered the calls to prayer, was that an empty ritual?

\hfill

\texttt{Proud Indo Hindu on 2015-12-19 at 20:14 said:}

``Reason attributes to the divine an attitude, a weakness, and even more often, purely human motives... which tended to humanize the divinities. The Good News announced by Jesus reversed this ancient conception, calling for the divinization of the human in man by a second Birth; the gateway to the Kingdom of God. ''

The beauty of Hinduism is that it does both. We humanize our gods and divinize ourselves.

\hfill

\texttt{infowarrior1 on 2015-12-23 at 22:11 said:}

The vikings suffered from grrl power syndrome thought at at least the TV series anyway

\hfill

\texttt{White Rabbit on 2015-12-24 at 01:22 said:}

Viking society was still patriarchal. Actually, women tend to be the weakest when they are ``free'':
\url{https://www.psychologytoday.com/blog/sexual-personalities/201504/are-women-more-emotional-men}


The fairer sex is strongest on its knees (a ``pretty picture'', in Trump's word).

\hfill

\texttt{Alistair Fraser on 2015-12-24 at 08:48 said:}

``Reason attributes to the divine an attitude, a weakness, and even more often, purely human motives... which tended to humanize the divinities. The Good News announced by Jesus reversed this ancient conception, calling for the divinization of the human in man by a second Birth; the gateway to the Kingdom of God.'

My interpretation of the second birth with regard to ``modern western life'' because that is what i am recently experiencing ... . in a simplified manner, is that so-called civilised citizenship existence in general ``pretends'' to be dictated by reason , and there are rules of law that also proclaim to adhere to an ``agreed reason''

Now , it seems there are many people that can get through their individual lives , in these western societies conforming with and to these laid down reasons for existence which in a simplified manner are in their optimum , get schooled , get further universalised schooled, get a career, get a wife or husband, get a nice house, even a car, have holidays, friends , and you may then even live a very happy existence or at least qualify for the conditions of a very happy existence laid down by this reasonings.

Anyone who follows the above protocol as dictated to by their society will then be judged by others and their very selfs on just how happy they are based on their acquirement of the above acquisitions . If they don't have the collection , they can reason that they are not quite happy enough or at all because of that lack .

Some others , may even, on that journey of societal reasoned development , start to feel that something else is not quite being satisfied , and to satisfy that, they might start going or go to the weekly church, which provides a sanctuary that does at least address this little x-factor in their consciousness that was not fully happy or even in a little distress , and in that church they will mix with other like minded people whom will either be attending for reasons of social status, or reasons of intangible needs, or they might be one of the few higher developed beings that are there to guide and nourish others .

The second birth -- As each person is sort of launched into earthly existence through the trauma of birth (modern western birth is industrialised -RdLaing) they sort of come to consciousness with a certain amount of propulsion dictated by their keepers and the society structure of schooling, and as they get older and more involved and investing in this societal/keeper propulsion , it follows that their consciousness has a central stream of more powerful influences that have been induced into them ... ..for their own good ... .is the reasoning , to help them survive and prosper in the society.

Now in some people , there already is , or begins to manifest in time, another deeper mode of consciousness , which is potentially in all people, but in some, this inner voice becomes louder and its nature is firstly of protest , and that protest is that , the propulsion that one has been induced by into society , is now becoming a drain on the vitality and this then leads to various reactions, some will still have their ambitious eye completely focussed on the ladder of acquisitions that they are on and these people will ignore or apply a plaster to that inner protest , like alcohol or drugs or pharma.

The first stage of repression or denial which can also develop into more severe stages.

But others will take this inner protest seriously and this is the second birth as what they do, is they realise that up until this moment, they have been a strange kind of passenger in their own existence with may or may not have gone the way they wanted or thought they wanted , but from this moment on, they come to realise that this inner protest or (conscience) might just be an interface to higher intelligence or a god aspect and this is the awakening of the divine possibility in they. From that moment in time, they come to recognise the possibility that there may be a divine guide that speaks through inner conscience but that also this may arrive through others inner consciences and so now the world of communications and existence takes on a whole new meaning and purpose . They now start to take control over their own existence , not to satisfy societies reasoned ambition although that may very well co-ordinate for some, their main focus is to arrange their lifestyle into a nourishing environment that they then can exist in with a portion of self-control that is propelled by this very inner voice . and in this world of propulsion , nothing is absolutely certain, destiny may change route at any moment , but that inner voice is no longer a cry of protest, it starts to develop into a voice of wisdom developed in its depth and effect by the selfs aptitude

But there is also great danger possible because when it is ascertained that there is a higher intelligence or divinity then also there must be a lower order iof influences and these lower orders of influences can also seem to penetrate and occupy the ``false conscience'' like a trojan horse .

So to cut a long story short, happiness , whether by societal reasoning or divine recognition and attempted development , is a path of obstacles for all

\hfill

\texttt{Proud Indo Hindu on 2015-12-26 at 23:13 said:}

@White Rabbit ``The fairer sex is strongest on its knees''

I'll assume you mean in prayer. Oral sex is a perversion that the West has mainstreamed and tried to export only after it adopted it from gay porn. You can ask both of your grandmothers if this was expected of them. That this has become a staple of the western erotic repetoire only shows how pornified you've become. And if a presidential candidate openly promotes it, only shows the level of your political discourse.

\hfill

\texttt{Proud Indo Hindu on 2015-12-26 at 23:15 said:}

@Avery, ``I met a Chinese guy who has adopted an ironic faith in Buddhism. When he sees a Buddha image, he puts his hands together, bows, gets down on both knees and prostrates himself in prayer. Except that he obviously didn't learn this behavior from his parents; he's urban elite and extremely overeducated. ''

What makes you think his parents didn't teach him this? As an elite and overeducated Hindu myself, I can assure you that we do pass on our spiritual traditions and rituals, with much cultural pride as well.

\hfill

\texttt{White Rabbit on 2015-12-27 at 01:21 said:}

Hindu: thank you for giving me the opportunity to clarify, apparently that was an unfortunate choice of words (I didn't mean to imply anything pornographic). The point is, in a natural and healthy order wives will venerate their husbands.

\hfill

\texttt{Proud Indo Hindu on 2015-12-27 at 21:01 said:}

``The point is, in a natural and healthy order wives will venerate their husbands.''

In traditional Hindu culture the wife is seen as an embodiment of the Divine Feminine/Goddess and the husband as the Divine Masculine/God. The two are to venerate each other. We have rituals that invoke this sense of sacredness in the human.

\hfill

\texttt{Alistair Fraser on 2015-12-28 at 04:08 said:}

Below the horizon = of the flesh = profane = falling = dominated and ruled by past time

Above the horizon = Transumanar = of the imaginary = ascending = dominated and ruled by the future time

On the Horizon = the perfect conjunction of the imaginary and the flesh = creates a separation from time =''The two are to venerate each other''

``We have rituals that invoke this sense of sacredness in the human''

What are these ? Who has achieved such ? maybe not possible in todays moments, so new possibles to be achieved .

\hfill

\texttt{Proud Indo Hindu on 2015-12-28 at 23:14 said:}

Alistair, in Hindu culture we have tons of holy days, festivals, and we celebrate them in pomp. Our weddings are several days long and involve rituals for both bride and groom. We also have festivals for honoring the spouse. Anyway, we've got tons of festivals in which we honor goddesses, gods, family members, saints and of course cows

Come to India, you'll see. Or hang out at a Hindu temple in your own country to get a glimpse.

\hfill

\texttt{Alistair Fraser on 2015-12-29 at 10:07 said:}

@PiH ``we have tons of ... . '' i suppose im trying to extract information about the extent of this so called ``sacredness'' and ``veneration'' in those relationships , and im a little sceptical about many festivals etc as they often are hijacked by the socialising perogative which in my view tends to dip the dynamics of relating below the horizon , although i am always ready to be pleasantly surprised by existence and others in it.

But this ``veneration'' and ``rituals to amplify its bandwidth and depth '' sounds like a decent protocol for at least developing a more profound and respectful form of relating that should not stagnate and shrink as a victim of time more expand and enrichen as a master of time .

But i have seen no evidence of such relationships in a western capitalist society.

Usually any longevity of intimate or other lived-in relating is a result of mutual benefits in the material and psychological security (or insecuritys) sense and the initial intimacy has evolved through variable stages from lust to respect to moments of contempt and worse . And these 3 varying stages in their temporal manifestations then serve as a cyclic mirror of the internal harmony or disharmony of the relationship

And somewhere in there , there may be a phenomena called ``love'' applied and used in the rhetoric

This western mythical version of ``love'' which by the way is probably never actually elicitated in elaborative words, then serves as the bottom line glue that welds this relating in an intuitive sense rather than a sensibilised elaborated sense which then when re-affirmed , acts as a fuel to oversome the down cycles of the 3 stages.

So the mythical love uses a mythical faith as a real power to control lesser moments of human relationship expressions

I use the term mythical for love and faith because their historical and present power of actualisation far outweighs any form of universal elaborated meaning given for them , a good example of the power of mythological idea right there used through an apparent shared intuition either from wisdom or even fear in some , the fear of loss.

What exactly does ``i love you'' mean for many ? and how has that term been abused by those using it that have no comparison to others use of it , and thats how , in the end it comes down to an intuition from one to another, if that term is used authentically or not so. And even then, their shared moment of intuited ``love'' agreement is but a small splash in the bigger picture of possible developments and then it can be arrived at, the great piece of wisdomn, that ignorance is indeed bliss

\hfill

\texttt{Olavsson on 2016-01-02 at 13:12 said:}

@Indo Hindu:

Regarding your claims about how apparently spiritual your modern Hindus are, even the absurdly Westernized and profanely over-educated ``elite'', what is your comment to Rama Coomaraswamy's statement that almost nobody have the time to practice a true Hindu dharma seriously any longer? Was his insight into the matter insufficient? I'm afraid you have a superficial understanding of what this really means. Having a cultural-religious identity and carrying on certain customs is hardly breathtakingly impressive. Such options are available to westerners as well. It doesn't make the civilization truly Dharmic in any complete sense, however.

\hfill

\texttt{Mark Citadel on 2016-01-04 at 08:08 said:}

I couldn't help but find a link between this production and the Trump campaign. The advertising for it apparently almost looked like a Trump ad until the `sieg heil!'. Might be coincidence, but the premise would seem to bolster America's weariness of far right ideas with the specter of Nazism. I agree that it looked a little plastic from the trailer.

\hfill

\texttt{1234kid on 2016-01-23 at 17:06 said:}

``I suspect not, since there is no gay character to explain to the Vikings how to act more manly.'' haha, zing to donovan, right?!

i guess the vikings could be doing some mannerbund thing off screen but still procreating and doing family stuff with the womenfolk.

\hfill

\texttt{Akira Kalashnikov on 2016-01-24 at 18:23 said:}

Those trying to understand the more sinister aspects of Donald Trump would be well advised to read ``The Three stigmata of Palmer Eldritch'' by PKD.

\end{sffamily}
\end{footnotesize}

